%% ===================================================================
\section{Methods}
\label{sec:methods}

Figure~\ref{fig:pipeline} summarises the pipeline. Every stage is implemented
in the accompanying repository and executed by a single command; every
number reported in Sections~\ref{sec:results}--\ref{sec:discussion} is
generated from its outputs.

\begin{figure*}[!t]
\centering
% Fig. 1 -- processing pipeline. Pure TikZ, monochrome, explicit grid.
\begin{tikzpicture}[
  font=\footnotesize,
  box/.style={draw, rectangle, minimum height=11mm, minimum width=30mm,
              text width=28mm, align=center, inner sep=1.5pt, line width=0.5pt},
  io/.style={box, fill=black!7},
  res/.style={box, line width=1pt},
  arr/.style={-{Stealth[length=1.8mm,width=1.5mm]}, line width=0.5pt},
  band/.style={font=\scriptsize\bfseries, anchor=west, inner sep=0pt}]
\def\dx{35mm}\def\dy{-25mm}
% band A: corpus generation
\node[io]  (osm)   at (0*\dx, 0)      {Building footprints\\(Overture Maps)};
\node[box] (scene) at (1*\dx, 0)      {Geo-referenced scene\\LoD1, $\sceneSize\times\sceneSize$\,km};
\node[box] (rt)    at (2*\dx, 0)      {Sionna RT, \nSitesPerCity{} sites\\\freqGHz{} and \freqTwoGHz\,GHz};
\node[io]  (links) at (3*\dx, 0)      {Path-loss links\\\cellSize\,m cells, $\le\maxPl$\,dB};
% band B: tiles
\node[box] (tile)  at (0*\dx, 1*\dy)  {Tiling\\$G_s=\tileSize$\,m};
\node[box, minimum height=8.5mm] (desc)  at (1*\dx, 1*\dy+5.2mm) {Descriptor $\mathbf{v}_t\in\mathbb{R}^{32}$};
\node[box, minimum height=8.5mm] (fit)   at (1*\dx, 1*\dy-5.2mm) {Labels $\boldsymbol{\theta}_t=(\gamma_t,P_{L_0,t})$};
\node[io]  (table) at (2*\dx, 1*\dy)  {Tile table\\\nCities{} cities};
% band C: transfer and evaluation
\node[box] (op)    at (0*\dx, 2*\dy)  {Transfer operator\\(source cities)};
\node[box] (fs)    at (1*\dx, 2*\dy)  {Few-shot correction\\($k$ surveyed tiles)};
\node[res] (pm)    at (2*\dx, 2*\dy)  {Parameter error\\$\gamma$, $P_{L_0}$};
\node[res] (net)   at (3*\dx, 2*\dy)  {Network metrics,\\planning, certificate};
\node[box] (rtn) at (3*\dx, 1*\dy) {Site-aware link\\transfer model};
% flows
\draw[arr] (osm) -- (scene);
\draw[arr] (scene) -- (rt);
\draw[arr] (rt) -- (links);
\draw[arr] ([xshift=-6mm]links.south) -- ++(0,-6.5mm) -| ([xshift=6mm]tile.north);
\draw[arr] (scene.south) -- (desc.north);
\draw[arr] (tile.east) -- ++(1.6mm,0) |- (desc.west);
\draw[arr] (tile.east) -- ++(1.6mm,0) |- (fit.west);
\draw[arr] (desc.east) -- ++(1.6mm,0) |- (table.west);
\draw[arr] (fit.east) -- ++(1.6mm,0) |- (table.west);
\draw[arr] (table.south) -- ++(0,-6.5mm) -| ([xshift=6mm]op.north);
\draw[arr] (op) -- (fs);
\draw[arr] (fs) -- (pm);
\draw[arr] (pm) -- (net);
\draw[arr] ([xshift=8mm]links.south) -- ([xshift=8mm]rtn.north);
\draw[arr] (rtn) -- (net);
% band labels (rotated, left margin)
\foreach \n/\t in {osm/(a) Corpus, tile/(b) Tiles, op/(c) Transfer}
  \node[band, rotate=90, anchor=south] at ([xshift=-2.5mm]\n.west) {\t};
\end{tikzpicture}

\caption{Processing pipeline. (a)~Geo-referenced scenes are built from
open building footprints and ray traced. (b)~Links are grouped into tiles; each tile
receives a label $\btheta_t=(\gamma_t,P_{L_0,t})$ and a descriptor
$\bv_t$ computed from the same scene geometry. (c)~Tile transfer operators
fitted on source cities predict the labels of a held-out city, optionally
corrected with $k$ surveyed tiles; the site-aware link model predicts every
link from its geometry. Both are evaluated through the network quantities of
Section~\ref{sec:netmodel}, a site plan and a coverage certificate.}
\label{fig:pipeline}
\end{figure*}

\subsection{Ray-traced corpus}
\label{sec:corpus}

Each city is a square of side \SI{\sceneSize}{\kilo\metre} centred on a
fixed coordinate and expressed in a local metric frame, the UTM projection
of its zone shifted to the centre~\cite{snyder1987}. Building footprints are
taken from the buildings theme of Overture Maps~\cite{overture}, release
\overtureRelease{},
a versioned open dataset whose primary source is
OpenStreetMap~\cite{haklay2008osm}; a pinned release makes the geometry
exactly reproducible, which a live map query is not. Because grid north and
true north differ, the footprint query box covers all four corners of the
square. Building heights are resolved in order from the height attribute,
from the number of floors at \SI{3.2}{\metre} per floor, and otherwise from
the median of the resolved heights of the city; the share of each source is
reported per city. Footprint exteriors are clipped to the square, parts
smaller than \SI{4}{\metre\squared} are discarded, and the remainder is
extruded to level-of-detail~1 and placed on a flat ground plane. Buildings
use the ITU concrete model and the ground the ITU medium-dry-ground
model~\cite{itup2040}.

Each city has \nSitesPerCity{} sites on rooftops: local height maxima of
the building mesh on a \SI{25}{\metre} grid are ranked by height and
selected greedily with a minimum separation of \SI{\minSep}{\metre}, and the
antenna is mounted \SI{\siteMast}{\metre} above the roof; if too few
candidates satisfy the separation, the remaining sites are drawn at random
from the other rooftop candidates with a fixed seed. Radio maps are
computed with Sionna~RT~\cite{hoydis2023sionna,hoydis2024learning}
(version~\sionnaVersion) at \SI{\freqGHz}{\giga\hertz} with isotropic
antennas, \rtSamples{} rays per site, maximum interaction depth
\rtDepth{}, specular and diffuse reflection, refraction and diffraction
(including edge diffraction) enabled, on a
\SI{\cellSize}{\metre} grid at \SI{\ueHeight}{\metre} receiver height. Cells
with path loss above \SI{\maxPl}{\decibel}, beyond any practical coupling
loss, are discarded. The solver version and settings are stored with every
city. \nSitesWord{} sites on \SI{\sceneArea}{\kilo\metre\squared} give one site per
\SI{\areaPerSite}{\kilo\metre\squared}, the density of a hexagonal layout with an
inter-site distance of about \SI{\isdMain}{\metre}, sparser than the
\SI{500}{\metre} of the 3GPP urban-macro scenario~\cite{tr38901}; ranking
rooftops by height follows macro-cell practice. Both choices are tested in
Section~\ref{sec:res-robust}.

Two further corpora use the same scenes and settings. The \emph{second
band} re-traces every site at \SI{\freqTwoGHz}{\giga\hertz}, in the
\SIrange{7}{24}{\giga\hertz} range studied for NR~\cite{tr38820}. The
\emph{siting replicate} re-traces \nRandCities{} cities (\randCities) at
\SI{\freqGHz}{\giga\hertz} with the same number of sites taken from the
rooftop candidates in random order instead of by height, under the same
separation rule.

\paragraph{Consistency check}
The planar radio-map estimator of the solver weights a ray crossing the
measurement plane by the inverse cosine of its incidence angle, and a
diffracted ray by an integration weight over the diffraction cone; the
first weight is unbounded and the second can be very large, so a cell
occasionally receives one overweighted contribution. Its path loss can then fall below free space by tens of
decibels, which no incoherent sum of a few propagation paths allows. We
remove every link whose path loss lies more than
\SI{\qcMargin}{\decibel} below the free-space loss at its 3-D distance:
\qcRemoved{} links, \qcShare\% of the corpus. For the site with the most
such links, disabling diffraction reduces the links more than
\SI{\qcMargin}{\decibel} below free space from \outDiffOn{} to \outDiffOff{}
with otherwise identical settings, which identifies the diffraction weights
as their source.

\subsection{Tile labels}
\label{sec:labels}

Links are assigned to tiles of side $G_s=\SI{\tileSize}{\metre}$ by their
receiver position; tiles with fewer than \minLinks{} links are not labelled.
Inside a tile, the ordinary least-squares (OLS) slope is poorly determined
because every receiver sees the serving sites from a similar distance.

\begin{proposition}[Ill-conditioning of the per-tile fit]
\label{prop:identify}
For $\mathbf{A}=[\mathbf{x}\;\mathbf{1}]$ with $n$ rows,
$\det(\mathbf{A}^{\top}\mathbf{A})=n^{2}\Var(x)$ and
$\Var(\hat\gamma_{\mathrm{OLS}})=\sigma_\varepsilon^{2}/(n\Var(x))$, which
diverges as the intra-tile spread of $x$ vanishes.
\end{proposition}

\noindent A receiver at distance $d$ from its site sees a tile of side $G_s$
span only $\Delta x\approx10\log_{10}(1+G_s/d)$, which is
\SI{0.41}{\decibel} for $G_s=\SI{100}{\metre}$ and $d=\SI{1}{\kilo\metre}$.
We therefore shrink the slope towards the city-wide mean with weights set by
the data (empirical Bayes~\cite{morris1983}).

\begin{proposition}[Empirical-Bayes shrinkage]
\label{prop:eb}
Let $b_t$ be the OLS slope of tile $t$, computed with $x$ centred inside
the tile, with sampling variance $v_t$, and let the true slopes vary between
tiles with mean $\mu$ and variance $\tau^{2}$. Among
estimators $\hat\gamma_t=(1-w)b_t+w\mu$ the mean squared error
$(1-w)^{2}v_t+w^{2}\tau^{2}$ is minimised by
\begin{equation}
w_t=\frac{v_t}{v_t+\tau^{2}}=\frac{\lambda_t}{\lambda_t+S_{xx,t}},
\qquad \lambda_t=\frac{S_{xx,t}\,v_t}{\tau^{2}},
\label{eq:eb}
\end{equation}
which is ridge regression on the centred slope with penalty $\lambda_t$
and prior mean $\mu$. Its risk is $v_t\tau^{2}/(v_t+\tau^{2})\le\min(v_t,\tau^{2})$,
bounded by $\tau^{2}$ however degenerate the tile geometry.
\end{proposition}

\noindent Radio-map cells served by the same site share propagation paths,
so the independent-residual variance $s_t^{2}/S_{xx,t}$ understates $v_t$.
We therefore use the cluster-robust variance of Eq.~\eqref{eq:cr} (below)
of the OLS slope, with the serving sites as clusters, never below the
independent value; a tile served by a single site uses the independent
value multiplied by the median ratio of the two over the city (its design
effect). We estimate $(\mu,\tau^{2})$ per city with the
DerSimonian--Laird moment estimator~\cite{dersimonian1986} over the tiles
with a non-degenerate span, set
$\hat\gamma_t=(1-w_t)b_t+w_t\hat\mu$, and leave the intercept unpenalised:
$\hat P_{L_0,t}=\bar y_t-\hat\gamma_t\bar x_t$. A tile with no span takes
$\hat\gamma_t=\hat\mu$. The shrinkage weight $w_t$ is published with every
label. Algorithm~\ref{alg:labels} states the procedure; it is vectorised
over all links of a city.

\begin{algorithm}[!t]
\SetAlgoLined\DontPrintSemicolon
\KwIn{links $(x_i,y_i,t_i,b_i)$ of one city; minimum count $n_{\min}$}
\KwOut{$\hat\btheta_t$, $w_t$, iid and cluster-robust standard errors}
\ForEach{tile $t$ with $n_t\ge n_{\min}$}{
  $\bar x_t,\bar y_t,S_{xx,t},S_{xy,t}$;\quad $b_t\leftarrow S_{xy,t}/S_{xx,t}$;\quad
  $v_t\leftarrow$ cluster-robust variance of $b_t$ (Eq.~\eqref{eq:cr})\;
}
$(\hat\mu,\hat\tau^{2})\leftarrow$ DerSimonian--Laird over $\{(b_t,v_t)\}$\;
\ForEach{tile $t$}{
  $w_t\leftarrow v_t/(v_t+\hat\tau^{2})$;\quad
  $\hat\gamma_t\leftarrow(1-w_t)b_t+w_t\hat\mu$;\quad
  $\hat P_{L_0,t}\leftarrow\bar y_t-\hat\gamma_t\bar x_t$\;
  residuals $e_i\leftarrow y_i-\hat P_{L_0,t}-\hat\gamma_t x_i$\;
  standard errors from Eq.~\eqref{eq:cr} with iid and with site clusters\;
}
\caption{Tile labels}
\label{alg:labels}
\end{algorithm}

\paragraph{Label precision}
The estimator is linear in the path losses. In the centred parametrisation
$y_i=a_t+\gamma_t(x_i-\bar x_t)+e_i$ with design rows
$\mathbf{z}_i=(1,x_i-\bar x_t)$, the coefficients are
$\mathbf{M}\sum_i\mathbf{z}_i y_i$ up to the prior term, with
$\mathbf{M}=\operatorname{diag}\bigl(1/n,\,1/(S_{xx,t}+\lambda_t)\bigr)$
($\lambda_t=0$ for the OLS slope), and $\hat P_{L_0,t}=\hat a_t-\hat\gamma_t\bar x_t$.
With the links of the tile partitioned by serving site into $G$ clusters
$\mathcal{I}_1,\dots,\mathcal{I}_G$, with design blocks $\mathbf{Z}_c$ and
residuals $\mathbf{e}_c$, the cluster-robust variance is
\begin{equation}
\widehat{\Var}_{\mathrm{cl}}=\mathbf{M}\Bigl(\sum_{c=1}^{G}\mathbf{u}_c\mathbf{u}_c^{\top}\Bigr)\mathbf{M},
\quad
\mathbf{u}_c=\mathbf{Z}_c^{\top}(\mathbf{I}-\mathbf{H}_{cc})^{-1/2}\mathbf{e}_c,
\label{eq:cr}
\end{equation}
with $\mathbf{H}_{cc}=\mathbf{Z}_c\mathbf{M}\mathbf{Z}_c^{\top}$ (Moore--Penrose
inverse root where $\mathbf{I}-\mathbf{H}_{cc}$ is singular). This is the
bias-reduced linearisation (CR2) of~\cite{bell2002,pustejovsky2018}; without
the factor $(\mathbf{I}-\mathbf{H}_{cc})^{-1/2}$ it is the usual
estimator~\cite{liang1986,cameron2015}, which is biased downwards when the
clusters are few, as here. Since
$\mathbf{Z}_c^{\top}f(\mathbf{Z}_c\mathbf{M}\mathbf{Z}_c^{\top})=f(\mathbf{Z}_c^{\top}\mathbf{Z}_c\mathbf{M})\mathbf{Z}_c^{\top}$
for $f(s)=(1-s)^{-1/2}$, every $\mathbf{u}_c$ requires only the eigen-decomposition
of a $2\times2$ matrix. The iid variance is
$s^{2}\,\mathbf{M}(\sum_i\mathbf{z}_i\mathbf{z}_i^{\top})\mathbf{M}$.
Because the number of clusters per tile is small,
Eq.~\eqref{eq:cr} is complemented by a split-half estimate: the sites of a
city are divided at random into two disjoint halves, the complete
estimator (including its own $\hat\mu,\hat\tau^2$) is refitted on each half,
and $\operatorname{sd}(\hat\gamma^{(1)}-\hat\gamma^{(2)})/\sqrt2$ estimates the
standard error of a half-data label; dividing by a further $\sqrt2$
extrapolates to the full data under the assumption that the variance scales
with the inverse number of sites. Five random partitions are pooled.

\subsection{Tile descriptors}
\label{sec:descriptors}

The descriptor must describe exactly the geometry the ray tracer saw. It is
therefore computed from the same footprint parts, heights and scene clipping
as the extruded mesh, with every part clipped to the tile it falls in. It
contains the $32$ built-environment features of~\cite{ozyurt2026}: footprint
ratio and building count in five height classes ($<5$, $5$--$10$,
$10$--$30$, $30$--$60$, $\ge\SI{60}{\metre}$); mean and standard deviation of
footprint area, height and volume; total footprint density; and, for
horizontal planes at $5$, $10$, $30$, $60$ and \SI{100}{\metre}, the number
of buildings reaching the plane and the mean and standard deviation of the
distance between their centroids. The topographic, land-cover and climatic
descriptors of~\cite{ozyurt2026} are omitted on purpose: the scenes have a
flat ground, no water or vegetation surfaces and no weather, so these
features would describe nothing the solver modelled and could only act as
city identifiers. Descriptors are stored unnormalised; every experiment
scales them by the column-wise maximum absolute value of its training
cities, never of the target.

\subsection{Transfer operators}
\label{sec:operators}

We compare constant predictors, label-averaging operators and learned
regressors.
\begin{itemize}
\item \emph{Source median}: the median source label, the no-information
reference.
\item \emph{3GPP UMa}: the line-of-sight and non-line-of-sight single-slope
formulas of TR~38.901~\cite{tr38901} at the reference distance, which give
fixed parameters for every tile, and \emph{COST-231 Hata}~\cite{cost231}
(metropolitan correction, median site height of the source cities), used
outside its validity range of \SIrange{1.5}{2}{\giga\hertz} and
\SIrange{1}{20}{\kilo\metre} as a further standard model.
\item \emph{Kernel (Nadaraya--Watson)}~\cite{nadaraya1964,watson1964}, the
operator of~\cite{ozyurt2026}:
\begin{equation}
\hat\btheta_u=\sum_j\tilde w_j(u)\btheta_j,\qquad
\tilde w_j(u)\propto\exp\!\Bigl(-\frac{\lVert\bv_u-\bv_j\rVert^{2}}{2\sigma^{2}}\Bigr),
\label{eq:nw}
\end{equation}
computed with the minimum squared distance subtracted in the exponent, which
leaves the weights unchanged and avoids numerical underflow.
\item \emph{GP-ARD}: a Gaussian process with one length scale per
descriptor (automatic relevance determination)~\cite{rasmussen2006}, i.e.\ a
kernel operator whose metric is learned; the hyper-parameters maximise the
marginal likelihood on \num{600} random source tiles (it needs no validation
city) and the posterior mean uses every source tile.
\item \emph{$k$-NN}: the mean label of the $k$ nearest source tiles.
\item \emph{Ridge}, a two-layer \emph{MLP}, gradient-boosted trees
(\emph{GBDT}, LightGBM~\cite{ke2017lightgbm}) and a \emph{neural encoder}
with a constrained $(\gamma,P_{L_0})$ head trained with a heteroscedastic
Gaussian likelihood.
\item Two domain-adaptation variants that use the unlabelled descriptors of
the target city: \emph{GBDT-IW} weights the source tiles by the density
ratio of target to source descriptors estimated with a logistic domain
classifier, clipped at its 99th percentile~\cite{shimodaira2000}, and
\emph{CORAL} re-colours the source descriptors with the target covariance
before ridge regression~\cite{sun2016coral}.
\end{itemize}
Each operator's hyper-parameters (bandwidth $\sigma$, $k$, ridge, CORAL and
MLP penalties, number of trees, early-stopping epoch) are chosen on the
validation city of the split by the root-mean-square standardised error over
both parameters; the chosen values are listed in
Table~\ref{tab:hparams}. Stochastic operators are run with five seeds.

\paragraph{Behaviour of kernel transfer}
Write $d_j=\lVert\bv_u-\bv_j\rVert^{2}$, $\delta_j=d_j-\min_k d_k\in[0,D_u]$,
$\bar\btheta=N^{-1}\sum_j\btheta_j$ and
$\sigma_\theta^{2}=N^{-1}\sum_j\lVert\btheta_j-\bar\btheta\rVert^{2}$.

\begin{theorem}[Two regimes of kernel transfer]
\label{thm:regimes}
(a)~If $\epsilon_u=D_u/(2\sigma^{2})<1$, then
\begin{equation}
\hat\btheta_u=\bar\btheta-\frac{1}{2\sigma^{2}}\Cov_j(\delta_j,\btheta_j)
+O(\epsilon_u^{2}),\quad
\bigl\lVert\hat\btheta_u-\bar\btheta\bigr\rVert\le
\frac{D_u\,\sigma_\theta}{4\sigma^{2}}+O(\epsilon_u^{2}),
\label{eq:flat}
\end{equation}
and $n_{\mathrm{eff}}(u)=N(1+O(\epsilon_u^{2}))$: the operator returns the
source mean up to a first-order correction.
(b)~Let $g_u$ be the gap between the smallest and second-smallest squared
distance. Then the weight of the nearest source tile satisfies
\begin{equation}
\tilde w_{(1)}\ge\bigl(1+(N-1)e^{-g_u/(2\sigma^{2})}\bigr)^{-1},
\label{eq:nn}
\end{equation}
so $\hat\btheta_u$ tends to the label of the nearest source tile when
$g_u\gg2\sigma^{2}\ln N$.
\end{theorem}

\noindent The minimum distance cancels from Eq.~\eqref{eq:nw}; a target
far from every source tile is therefore \emph{not} by itself pushed to
either regime. Which regime applies is decided by the selected bandwidth
relative to the spread $D_u$ and the gap $g_u$, and we measure it
(Section~\ref{sec:res-kernel}) rather than assume it. Both regimes carry a
failure mode under shift: the flat regime returns the source mean and
ignores the target descriptor; the nearest-neighbour regime copies one
source tile and transmits its label noise undamped
(Proposition~\ref{prop:decomp} with $n_{\mathrm{eff}}=1$).

\subsection{Few-shot correction}
\label{sec:fewshot}

With $k$ target tiles surveyed, the zero-shot prediction $g(\bv)$ is
corrected either by the mean anchor residual (\emph{bias}) or by a ridge
model of the residual on the descriptors (\emph{residual},
penalty $10$ on normalised descriptors). Anchors are drawn uniformly
at random, as the $k$ tiles closest to a random tile (\emph{contiguous}),
which mimics a single drive-test area~\cite{johansson2012mdt}, or by a
\emph{design} that uses no measurement: the $k$ tiles that greedily maximise
the facility-location objective
$\sum_{i}\max_{j\in T}\exp(-\lVert\bv_i-\bv_j\rVert^{2}/2h^{2})$ over the target
tiles, with $h$ the median pairwise descriptor distance; the objective is
monotone submodular, so the greedy set is within a factor $(1-1/e)$ of the
optimum~\cite{nemhauser1978,krause2008}. Evaluation is
always on the unsurveyed tiles, anchors are shared by all operators, and
ten random draws are made per $k\in\{3,5,10,25,50,100\}$.

\subsection{Site-aware link transfer}
\label{sec:link}

Proposition~\ref{prop:tilemodel} shows that any receiver-tile model assigns every
pixel to its nearest site and describes interference through the exponent
alone. As a transfer model free of this restriction we predict the path loss
of each (pixel, site) link directly. The features are computed from the
scene geometry only, so they are available for a city without any
measurement. The building heights are rasterised at \SI{2}{\metre} and the
straight segment from the site to the pixel is sampled at $96$ interior
points, where the building height $H(s)$ is compared with the
line-of-sight height $h_{\mathrm{los}}(s)$. The link features are the 3-D and
horizontal distances, the site height and the elevation angle; the fraction
of the segment with $H>h_{\mathrm{los}}$, the number of separate
obstructions and the length over footprints; the largest excess height
$\max_s(H-h_{\mathrm{los}})$; the single knife-edge loss
$J(\nu)=6.9+20\log_{10}\bigl(\sqrt{(\nu-0.1)^{2}+1}+\nu-0.1\bigr)$ for
$\nu>-0.78$ (zero otherwise) of the most obstructing point, with
$\nu=h\sqrt{2(1/d_1+1/d_2)/\lambda}$~\cite{itup526}; whether the receiver lies
inside a footprint and that building's height; and the built fraction and
mean building height in a \SI{100}{\metre} square window centred on the
receiver and the built fraction in a \SI{200}{\metre} window centred on the
site. The $32$ descriptors of the receiver's
tile are appended.

The ray tracer reports only links below the \SI{\maxPl}{\decibel} cut-off,
so a regressor trained on them never learns that a distant, obstructed site
is not received at all; it would predict a finite loss for every link and
overstate the reach of every site. We therefore use a hurdle model: a
gradient-boosted classifier~\cite{ke2017lightgbm} of whether a (pixel, site)
pair lies below the cut-off, trained on \num{120000} pairs of each training
city (all sites and the pixels of its labelled tiles, including the pairs
beyond the cut-off), and a gradient-boosted regressor
of the path loss, trained on \num{120000} ray-traced links of each training
city. A link is predicted beyond the cut-off (zero received power) when the
classifier gives a probability below one half, and otherwise its path loss
is the regression. Both models use early stopping on the validation city
(learning rate $0.1$, $127$ leaves, at most $1500$ trees, five seeds). The
path-loss error of the link model is that of the regressor over the
ray-traced links; the regressor alone, which keeps every link, is reported as
an ablation.

\paragraph{Survey design for site calibration}
In the few-shot variant the links observed in $k$ surveyed tiles, which carry
their serving site as minimisation-of-drive-tests (MDT) reports
do~\cite{johansson2012mdt,ts37320}, calibrate a per-site offset. Model the
residual of a link of site $s$ as $o_s+\varepsilon$, with independent site
offsets $o_s$ of prior mean $\bar o$ and variance $\tau_o^{2}$ and link noise of
variance $\sigma_\varepsilon^{2}$, and let $\kappa=\sigma_\varepsilon^{2}/\tau_o^{2}$.
With $n_s$ observed links of site $s$ and mean residual $\bar r_s$, the
posterior mean of $o_s$ is $(n_s\bar r_s+\kappa\bar o)/(n_s+\kappa)$, which we
use with $\bar o$ the mean residual of all observed links and
$\kappa=\kappaLink{}$ links, and its posterior variance is
$\tau_o^{2}\kappa/(n_s+\kappa)$.

\begin{proposition}[Survey design]
\label{prop:survey}
Let $c_{ts}\ge0$ be the number of links of site $s$ observable in tile $t$
and $n_s(T)=\sum_{t\in T}c_{ts}$ for a set $T$ of surveyed tiles. The total
reduction of posterior variance of the site offsets,
\begin{equation}
F(T)=\tau_o^{2}\sum_{s}\frac{n_s(T)}{n_s(T)+\kappa},
\label{eq:survey}
\end{equation}
is monotone and submodular in $T$, so the greedy choice of $k$ tiles attains
$F(T_{\mathrm{greedy}})\ge(1-1/e)\max_{|T|=k}F(T)$.
\end{proposition}

\noindent Before any measurement $c_{ts}$ is unknown; the \emph{designed}
survey uses the number of links of site $s$ in tile $t$ that the hurdle
classifier of the zero-shot link model predicts below the cut-off. It is compared
with random surveys for $k\in\{10,25,50\}$ (ten draws).

\subsection{Network-level evaluation}
\label{sec:neteval}

For the held-out city, each operator's tile parameters and the true site
positions define a predicted path-loss map for every (pixel, site) pair
through Eq.~\eqref{eq:logdist}; the ray-traced map is the truth, with pairs
beyond the \SI{\maxPl}{\decibel} cut-off treated as zero received power.
With $P_{\mathrm{tx}}=\SI{46}{dBm}$, \SI{20}{\mega\hertz} bandwidth and a
\SI{7}{\decibel} noise figure ($N=\SI{-94.0}{dBm}$) we compute the path-loss
error over all ray-traced links, the serving site, coverage at
$L\in\{120,130,140\}$\,dB, the SINR of Eq.~\eqref{eq:sinr} and the spectral
efficiency (attenuated Shannon, factor $0.6$, zero below \SI{-10}{\decibel},
capped at \SI{4.4}{b/s/Hz}~\cite{tr36942}). An SINR below
\SI{\sinrFloor}{\decibel}, including a pixel at which no site is predicted to
be received, is set to \SI{\sinrFloor}{\decibel} (outage) for the truth and
every predictor alike. Mobility decisions depend on
where the serving cell changes: a pixel is a cell-edge pixel when its
eastern or northern neighbour has a different serving site (one pixel per
boundary crossing), the \emph{edge
displacement} is the mean distance between the true and the predicted cell
edges (the average of the two directed mean nearest-edge distances), and the
\emph{edge-length error} is the relative error of the number of cell-edge
pixels, a proxy for the handover rate along a route.

\paragraph{Site planning}
Selecting $K$ of the $|\mathcal{B}|$ sites to maximise the number of covered
pixels is a monotone submodular maximisation, for which the greedy algorithm
attains a $(1-1/e)$ approximation of the optimum of the function it is
given~\cite{nemhauser1978}; site selection with propagation models is a
classical planning task~\cite{amaldi2003}. We run the greedy on the
\emph{predicted} coverage at \SI{140}{\decibel} and score the selected sites
on the \emph{true} coverage. The regret is the true coverage of the greedy
plan computed on the true map minus that of the plan computed on the
predicted map. For dimensioning, we also count the sites that the greedy
order on the predicted map needs to cover \covTarget\% of the pixels, and
the true coverage those sites achieve.

\paragraph{Certified coverage}
A planner must decide which pixels to \emph{promise} as covered. For a
transfer model and a held-out city $c$, let
$r_p=PL_{\mathrm{serv}}(p)-\widehat{PL}_{\mathrm{serv}}(p)$ be the true minus the
predicted serving path loss of pixel $p$ (the smallest path loss over the
sites), and $q_c$ the $(1-\alpha)$ quantile of $\{r_p\}$ over the pixels of
$c$.

\begin{proposition}[City-level coverage certificate]
\label{prop:certify}
Let the $n$ cities be exchangeable and the training procedure symmetric in
its source cities (the validation city drawn at random). With $q_c$ computed
from the leave-one-city-out model of each city $c$, let $\hat q_t$ be the
$j$-th largest of $\{q_c\}_{c\ne t}$ for a target city $t$, and claim pixel
$p$ as covered at level $L$ when $\widehat{PL}_{\mathrm{serv}}(p)+\hat q_t\le L$.
Then, with probability at least $(n-j)/n$ over the draw of the target city,
at most a fraction $\alpha$ of its pixels are claimed but not covered.
\end{proposition}

\noindent The guarantee is at the level of cities, the unit of shift, in the
spirit of conformal prediction~\cite{vovk2005,barber2021}; it needs no model
of the residuals. The rank $j$ trades confidence for coverage: $j=1$ (the
largest margin of the other cities) gives probability $(n-1)/n$, $j=2$ gives
$(n-2)/n$ with a smaller margin. We report it at $L=\SI{140}{\decibel}$ for
$\alpha\in\{0.05,0.10\}$ and $j\in\{1,2\}$ (the results at
\SI{130}{\decibel} are released with the code) with the claimed and the
falsely claimed share of the pixels, and compare it with the naive claim
$\widehat{PL}_{\mathrm{serv}}(p)\le L$.

\subsection{Robustness analyses}
\label{sec:robust}

\emph{Mixed-path composition.} The framework of~\cite{ozyurt2026}, following
MAPLE~\cite{ozyurt2020maple}, composes the loss of a link from every tile
its straight path crosses. We split the horizontal segment from each site to
each pixel at the tile borders (sampled every \SI{10}{\metre}), so that tile
$t$ holds a length $L_t$ of it, between the distances $d_t^{\mathrm{in}}$
and $d_t^{\mathrm{out}}$ from the site, and evaluate three compositions. The
first is the sum as written in the published framework,
$\PL=\sum_t[A_t+10\gamma_t\log_{10}L_t]$ with $L_t$ in metres, in which every
crossed tile adds its own log-distance loss. The second reads the same
algorithm as a multi-slope law along the path,
$\PL=A_{s}+\sum_t 10\gamma_t\log_{10}\big(\max(d_t^{\mathrm{out}},d_0)/
\max(d_t^{\mathrm{in}},d_0)\big)$, with the intercept of the site's tile $s$
and, over each stretch of distance, the exponent of the tile it lies in
(piecewise). The third averages the parameters along the path,
$\PL=\sum_t (L_t/d_h)\,[A_t+\gamma_t x]$ with the horizontal length $d_h$.
The piecewise and path-averaged compositions reduce to
Eq.~\eqref{eq:logdist} when the whole path lies in one tile. All three are
linear in the tile parameters, which are fitted jointly to all links of a
city by least squares with a weak ridge penalty ($\lambda=10$) towards the
best city-wide pair; they are evaluated with tiles of \SI{100}{\metre} and,
as in~\cite{ozyurt2026}, of \SI{200}{\metre}, \SI{500}{\metre} and
\SI{1000}{\metre}, always on the pixels of the main corpus. The labels of the tiles
crossed by at least 30 links are transferred with the source median, kernel
transfer and GBDT under the protocol of Section~\ref{sec:protocol} (one
seed), and the predicted maps are scored as in Section~\ref{sec:neteval}.
\emph{Network size.} Every zero-shot predictor is also scored on the network
of the first ten sites selected by the siting rule (one site per
\SI{0.9}{\kilo\metre\squared}, inter-site distance about
\SI{\isdSmall}{\metre}). \emph{Siting policy.} The models trained on the
main corpus are applied unchanged to the siting replicate; the tile model is
scored with labels fitted on the replicate itself and with the labels of the
same tiles from the main corpus. \emph{Tile size.} Labels, descriptors,
zero-shot transfer (median, ridge, GBDT) and the network evaluation are
repeated for $G_s\in\{50,100,200\}$\,m. \emph{Censoring.} Links beyond the
cut-off are missing, which can bias a least-squares slope. We refit the
labels with the cut lowered to $140$ and \SI{130}{\decibel} and, for
\num{300} random tiles per city, compare the least-squares slope with that of
a censored-Gaussian (Tobit) model~\cite{tobin1958} that also uses the
(pixel, site) pairs beyond the cut. \emph{Second band.} The tile and link
pipelines are repeated at \SI{\freqTwoGHz}{\giga\hertz}, and the
\SI{\freqGHz}{\giga\hertz} link models are applied to the second band with
the free-space offset $20\log_{10}(f_2/f_1)$, their knife-edge feature
evaluated at the new wavelength. \emph{External scenes.} The Munich, \'Etoile
and Florence scenes distributed with Sionna~RT were modelled independently
of our corpus, with detailed meshes and several ITU materials (brick,
concrete, marble, metal, wood). We derive building footprints and heights
from their meshes by vertical ray casting on a \SI{2}{\metre} grid, place
\nExtSites{} rooftop sites per scene with the rule of our corpus (minimum
separation \SI{\extMinSep}{\metre}, the scenes being smaller), ray trace them
with the corpus settings and evaluate, without any retraining, the tile
labels, the source median and GBDT fitted on all eleven cities, and the mean
of the eleven leave-one-city-out link models, zero-shot and with the per-site
calibration of Section~\ref{sec:link} from $k\in\{10,25\}$ surveyed tiles
(ten random draws, and the design of Proposition~\ref{prop:survey}). The San Francisco scene has
terrain and would require a terrain-following measurement surface; it is not
used.

\subsection{Protocol and inference}
\label{sec:protocol}

Every city is the target once (leave-one-city-out); one of the remaining
cities, drawn with a fixed seed, is the validation city and the others are
training cities. Tiles of one city are spatially correlated, so tile-level
tests would overstate significance; the unit of inference is the target city.
For each comparison we report the mean per-city difference, a 95\%
percentile bootstrap interval over cities~\cite{efron1993}, the exact
two-sided Wilcoxon signed-rank $p$-value~\cite{wilcoxon1945} and the number
of cities in which the first operator is better, with Holm's
correction~\cite{holm1979} within each family of comparisons. Standard
deviations across cities use $n-1$.
